\documentclass[11pt,a4paper]{article}

\usepackage[margin=2.2cm]{geometry}
\usepackage{amsmath,amssymb,amsthm}
\usepackage{booktabs}
\usepackage{array}
\usepackage{tabularx}
\usepackage[shortlabels]{enumitem}
\usepackage{xcolor}
\usepackage{tikz}
\usepackage{fancyhdr}
\usepackage{setspace}
\usepackage{microtype}
\usepackage{titlesec}
\usepackage{hyperref}
\usepackage{listings}
\usepackage[skins,breakable]{tcolorbox}

\usetikzlibrary{positioning,arrows.meta,calc,fit,shapes.geometric,matrix,decorations.pathreplacing}

\definecolor{ink}{HTML}{222222}
\definecolor{rulegray}{HTML}{444444}
\definecolor{boxline}{HTML}{666666}
\definecolor{boxfill}{HTML}{F4F4F4}
\definecolor{boxtitle}{HTML}{E8E8E8}
\definecolor{g1}{HTML}{D0D0D0}
\definecolor{g2}{HTML}{9A9A9A}
\definecolor{g3}{HTML}{5A5A5A}

\hypersetup{
  colorlinks=true,
  linkcolor=black,
  urlcolor=black,
  pdftitle={MTL 458 Practice Questions: Lectures 15--19},
  pdfauthor={Minati De, IIT Delhi}
}

\setstretch{1.12}
\raggedbottom
\setlist[itemize]{leftmargin=1.3em,itemsep=0.28em,topsep=0.35em}
\setlist[enumerate]{leftmargin=1.5em,itemsep=0.35em,topsep=0.4em}

\titleformat{\section}
  {\Large\bfseries\color{ink}}{\thesection.}{0.6em}{}
  [{\color{rulegray}\titlerule[0.5pt]}]
\titleformat{\subsection}
  {\large\bfseries\color{ink}}{\thesubsection}{0.55em}{}
\titleformat{\subsubsection}
  {\normalsize\bfseries\color{ink}}{\thesubsubsection}{0.5em}{}

\titlespacing*{\section}{0pt}{1.35em}{0.7em}
\titlespacing*{\subsection}{0pt}{1.05em}{0.45em}

\pagestyle{fancy}
\fancyhf{}
\renewcommand{\headrulewidth}{0.4pt}
\renewcommand{\footrulewidth}{0.3pt}
\fancyhead[L]{\small\color{ink}MTL 458 \textbullet\ Operating Systems}
\fancyhead[R]{\small\color{ink}Practice Questions: Lectures 15--19}
\fancyfoot[C]{\small\thepage}
\setlength{\headheight}{14pt}

% ---------------------------------------------------------------------------
% Boxes
% ---------------------------------------------------------------------------
\tcbset{
  commonbox/.style={
    colback=boxfill, colframe=boxline, boxrule=0.5pt,
    left=6pt, right=6pt, top=5pt, bottom=5pt,
    fonttitle=\bfseries, coltitle=black, colbacktitle=boxtitle,
    enhanced, sharp corners=south, breakable
  }
}
\newtcolorbox{keybox}[1][]{commonbox, title={#1}}
\newtcolorbox{quizbox}[1][]{commonbox, title={#1}}
\newtcolorbox{solbox}[1][]{commonbox, title={#1}}
\newtcolorbox{warnbox}[1][]{commonbox, title={#1}}
\newtcolorbox{defbox}[1][]{commonbox, title={#1}}

\lstset{
  language=C,
  basicstyle=\ttfamily\small,
  keywordstyle=\bfseries,
  commentstyle=\itshape\color{g3},
  numbers=left, numberstyle=\scriptsize\color{g2}, numbersep=6pt,
  frame=single, rulecolor=\color{boxline},
  backgroundcolor=\color{boxfill},
  xleftmargin=1.6em, framexleftmargin=1.4em,
  columns=fullflexible, keepspaces=true,
  showstringspaces=false, breaklines=true,
  aboveskip=0.6em, belowskip=0.6em,
  morekeywords={cond_t,mutex_t,sem_t,lock_t,rwlock_t,pthread_t,pthread_cond_t,pthread_mutex_t}
}
\lstdefinestyle{term}{language={}, numbers=none, basicstyle=\ttfamily\footnotesize,
  xleftmargin=0.4em, framexleftmargin=0.2em}

% Start a new page unless at least #1 of vertical space remains
\makeatletter
\newcommand{\keepspace}[1]{\par\penalty-100\begingroup
  \setlength{\dimen@}{#1}\dimen@ii\pagegoal\advance\dimen@ii-\pagetotal
  \ifdim\dimen@>\dimen@ii\ifdim\dimen@ii>\z@\vfil\fi\break\fi\endgroup}
\makeatother

\newcounter{qn}[section]
\newcommand{\lecno}{}
% \Q{type}{title}: numbered question box; solution follows in \begin{solbox}
\newenvironment{Qn}[2]{\refstepcounter{qn}\keepspace{8\baselineskip}%
  \begin{quizbox}[Q\lecno.\theqn\ \ [#1] #2]}{\end{quizbox}}
\newenvironment{Sol}{\begin{solbox}[Solution]}{\end{solbox}\vspace{0.4em}}

\begin{document}

\begin{center}
{\color{rulegray}\rule{\textwidth}{1.1pt}}\\[0.45em]
{\LARGE\bfseries Operating Systems}\\[0.25em]
{\Large\bfseries MTL 458 \textperiodcentered\ Practice Questions}\\[0.35em]
{\LARGE\bfseries Lectures 15--19}\\[0.2em]
{\large Swapping \textperiodcentered\ Concurrency \textperiodcentered\ Locks
\textperiodcentered\ Condition Variables \textperiodcentered\ Semaphores}\\[0.35em]
{\color{rulegray}\rule{\textwidth}{1.1pt}}
\end{center}

\vspace{0.4em}
\begin{keybox}[How to use this set]
\begin{itemize}
  \item Questions are grouped by lecture; each one is tagged with its type:
        \textbf{[T/F]} true/false with justification, \textbf{[Concept]},
        \textbf{[Trace]} simulate step by step, \textbf{[Bad schedule]}
        find an interleaving that breaks the code, \textbf{[Fill]} complete
        missing code/values, \textbf{[Predict]} possible outputs/results,
        \textbf{[Design]} write synchronization code.
  \item Within a lecture they go roughly from easy to hard. The style follows
        past MTL458 papers (fault tables, ticket-lock states, racy list
        outcomes, Peterson variants, semaphore design).
  \item Section 1 has all past-year exam questions on these lectures; the remaining sections are new practice questions.
  \item Try each question before reading its solution box. Everything uses
        only the content of Lectures 15--19 (and the notes).
\end{itemize}
\end{keybox}

\tableofcontents
\vspace{0.6em}
\noindent{\color{rulegray}\rule{\textwidth}{0.4pt}}

% ===========================================================================
\section{Past-year exam questions (Lectures 15--19)}
% ===========================================================================

\begin{keybox}[What this section contains]
Every question from the past MTL458 papers in \texttt{pyqs/} that falls in
Lectures 15--19 (page replacement, races, locks, condition variables,
semaphores), with full solutions. Questions on other lectures (paging
arithmetic, scheduling, deadlock detection/Banker's, process API, RAID)
are left out. The lefty/righty dining philosophers question appeared in
both Major 2016 and Minor~2 2016; it is included once.
\end{keybox}

\subsection*{Major 2016 --- Q1(c): LRU page replacement}
\paragraph{Restatement.}
Consider the following page reference strings:
\[
1,2,3,4,2,1,5,6,2,1,2,3,7,6,3,2,1,2,3,6.
\]
How many page faults would occur for the LRU replacement algorithm,
assuming four frames?

\paragraph{Solution.}
The four frame columns below denote fixed frame slots, not recency
order.  A dash in the victim column means either a hit or use of a
previously empty frame.
\[
\begin{array}{c|c|c|c|c|c|c|c}
\mathrm{Step}&\mathrm{Ref.}&F_1&F_2&F_3&F_4&\mathrm{Result}&\mathrm{Victim}\\ \hline
1 &1&1&-&-&-&\mathrm{fault}&-\\
2 &2&1&2&-&-&\mathrm{fault}&-\\
3 &3&1&2&3&-&\mathrm{fault}&-\\
4 &4&1&2&3&4&\mathrm{fault}&-\\
5 &2&1&2&3&4&\mathrm{hit}&-\\
6 &1&1&2&3&4&\mathrm{hit}&-\\
7 &5&1&2&5&4&\mathrm{fault}&3\\
8 &6&1&2&5&6&\mathrm{fault}&4\\
9 &2&1&2&5&6&\mathrm{hit}&-\\
10&1&1&2&5&6&\mathrm{hit}&-\\
11&2&1&2&5&6&\mathrm{hit}&-\\
12&3&1&2&3&6&\mathrm{fault}&5\\
13&7&1&2&3&7&\mathrm{fault}&6\\
14&6&6&2&3&7&\mathrm{fault}&1\\
15&3&6&2&3&7&\mathrm{hit}&-\\
16&2&6&2&3&7&\mathrm{hit}&-\\
17&1&6&2&3&1&\mathrm{fault}&7\\
18&2&6&2&3&1&\mathrm{hit}&-\\
19&3&6&2&3&1&\mathrm{hit}&-\\
20&6&6&2&3&1&\mathrm{hit}&-
\end{array}
\]
Faults occur at steps
\[
1,2,3,4,7,8,12,13,14,17,
\]
so the exact number of page faults is
\[
\boxed{10}.
\]

\subsection*{Major 2016 --- Q2(c) / Minor 2 2016 --- Q3: Lefty/righty dining philosophers}
\paragraph{Restatement.}
Suppose that there are two types of philosophers.  One type always
picks up his left fork first (a ``lefty''), and the other type always
picks up his right fork first (a ``righty'').  Using semaphore, define
both types of philosophers' behavior.

\paragraph{Solution.}
For \(n\) philosophers, let each fork be a binary semaphore:
\[
\mathrm{semaphore}\;\mathrm{fork}[0\ldots n-1],
\qquad \mathrm{fork}[k]=1\quad\hbox{initially}.
\]
For philosopher \(i\), define
\(\mathrm{left}(i)=i\) and
\(\mathrm{right}(i)=(i+1)\bmod n\).  Here \(P\) is wait/down and \(V\)
is signal/up.

\begin{verbatim}
lefty(i):
    while (true) {
        think();
        P(fork[left(i)]);
        P(fork[right(i)]);
        eat();
        V(fork[right(i)]);
        V(fork[left(i)]);
    }

righty(i):
    while (true) {
        think();
        P(fork[right(i)]);
        P(fork[left(i)]);
        eat();
        V(fork[left(i)]);
        V(fork[right(i)]);
    }
\end{verbatim}

Each fork is used by at most one philosopher because its semaphore is
binary.  If the circle contains both acquisition orientations (the
usual implementation makes one philosopher a righty and all the
others lefties), the uniform circular-wait pattern is broken, so all
philosophers cannot deadlock while holding their first fork.  If every
philosopher has the same orientation, the familiar circular deadlock
is still possible.

\subsection*{Major 2016 --- Q3(a): All concurrent traces}

\paragraph{Restatement.}
Consider a concurrent program with two processes, \(p\) and \(q\),
defined as follows.  \(A,B,C,D,\) and \(E\) are arbitrary atomic
(indivisible) statements.  Assume that the main program (not shown)
does a \texttt{parbegin} of the two processes.
\begin{verbatim}
Void p()             void q()
{ A; B; C; }         { D; E; }
\end{verbatim}
Show all the possible interleavings of the execution of the preceding
two processes (show this by giving execution ``traces'' in terms of the
atomic statements).

\paragraph{Solution.}
Every trace must preserve \(A<B<C\) and \(D<E\).  Choosing the three
positions occupied by \(A,B,C\) gives
\[
\binom{5}{3}=\binom{5}{2}=10
\]
traces.  All ten are
\[
\begin{array}{lllll}
ABCDE, & ABDCE, & ABDEC, & ADBCE, & ADBEC,\\
ADEBC, & DABCE, & DABEC, & DAEBC, & DEABC.
\end{array}
\]

\subsection*{Major 2016 --- Q3(b): Peterson's algorithm}

\paragraph{Restatement.}
Prove the correctness of Peterson's algorithm for mutual exclusion and
progress of some processes requiring access to its critical section.

\paragraph{Solution.}
The shared variables are initialized by
\[
\mathrm{flag}[0]=\mathrm{flag}[1]=\mathrm{false};
\]
\(\mathrm{turn}\) may initially have either value.  Process \(P_i\),
where \(j=1-i\), executes:

\begin{verbatim}
while (true) {
    flag[i] = true;
    turn = j;
    while (flag[j] && turn == j)
        ;                       /* spin */

    critical_section();

    flag[i] = false;
    remainder_section();
}
\end{verbatim}

This proof assumes atomic reads and writes, sequentially consistent
shared memory, and that a process does not remain in its critical
section forever.

\paragraph{Mutual exclusion.}
Assume both \(P_0\) and \(P_1\) are simultaneously in their critical
sections.  Both flags must then be true.  Consider the two entry-section
writes to \(\mathrm{turn}\), and suppose without loss of generality that
\(P_0\)'s write \(\mathrm{turn}=1\) occurs last.  Before that write,
\(P_1\) has already set \(\mathrm{flag}[1]\) to true, and that flag
cannot be cleared before \(P_1\) leaves its critical section.  Thus
\(P_0\) subsequently tests
\(\mathrm{flag}[1]\mathbin{\&\&}\mathrm{turn}=1\), which is true, and
cannot enter while \(P_1\) is in its critical section.  This
contradicts the assumption that both entered.  The case in which
\(P_1\)'s turn write is last is symmetric.  Therefore at most one
process is in the critical section.

\paragraph{Progress and bounded waiting.}
If only \(P_i\) wants to enter, then \(\mathrm{flag}[j]\) is false and
\(P_i\) enters immediately.  If both want to enter, both set their
flags, and the last write to \(\mathrm{turn}\) determines which process
waits.  Exactly one loop condition can remain true, so the other
process enters; there is no state in which both wait for each other.
On leaving, that process clears its flag, enabling the waiting process.
Moreover, after \(P_i\) announces its interest, \(P_j\) can enter ahead
of it at most once before \(P_i\) enters.  Thus the algorithm supplies
mutual exclusion, progress, and bounded waiting under fair execution.

\subsection*{Major 2016 --- Q3(c): An incorrect bounded-buffer implementation}

\paragraph{Restatement.}
Provide the algorithms for producer and consumer processes in bounded
buffer problem and show that they may not function correctly when
executed concurrently.

\paragraph{Solution.}
The following tempting implementation uses a circular buffer of size
\(N\) and a shared count, but no atomic operation, lock, or semaphore:

\begin{verbatim}
/* shared: item buffer[N]; int in = 0, out = 0, count = 0; */

produce(x):
    while (count == N)
        ;                       /* wait for a free slot */
    buffer[in] = x;
    in = (in + 1) % N;
    count = count + 1;

consume():
    while (count == 0)
        ;                       /* wait for an item */
    x = buffer[out];
    out = (out + 1) % N;
    count = count - 1;
    return x;
\end{verbatim}

The source-level increments are read--modify--write sequences.  For
example, suppose initially \(\mathrm{count}=1\), and a producer and a
consumer run concurrently:
\[
\begin{array}{c|l}
1&\mathrm{producer\ reads\ count}=1\\
2&\mathrm{consumer\ reads\ count}=1\\
3&\mathrm{producer\ computes}\ 2\\
4&\mathrm{consumer\ computes}\ 0\\
5&\mathrm{producer\ stores}\ 2\\
6&\mathrm{consumer\ stores}\ 0
\end{array}
\]
One item was added and one removed, so the correct final count is
\(1\), but the stored count is \(0\).  Reversing the final two stores
makes it \(2\), which is also wrong.  Future executions can therefore
underflow, overflow, or wait despite the actual buffer state.

With multiple producers, two producers can also both observe
\(\mathrm{count}=N-1\), pass the full test, use the same value of
\(\mathrm{in}\), and overwrite one another's item.  Analogous races
exist for multiple consumers at \(\mathrm{out}\).  Correct bounded
buffer code needs mutual exclusion around buffer/index updates and
counting synchronization, conventionally semaphores
\(\mathrm{empty}=N\), \(\mathrm{full}=0\), and
\(\mathrm{mutex}=1\).

\subsection*{Major 2020 --- Q1(c),(d): True or false, with justification}
\begin{enumerate}
\item[(c)] \emph{Threads \(T_1\) and \(T_2\) share \(x\).  \(T_1\) protects
\(x\) using a POSIX mutex.  If \(T_2\) writes \(x\) without locking the
mutex, a trap is generated.}

\textbf{False.}  A POSIX mutex is a software synchronization convention;
ordinary memory accesses do not carry information about which mutex should
protect them.  The hardware therefore does not trap merely because \(T_2\)
failed to lock the mutex.  The accesses constitute a data race (and hence
undefined behaviour in the C/POSIX memory model) if they are conflicting and
concurrent, but no protection trap is inherently generated.

\item[(d)] \emph{A thread cannot go to sleep while holding a lock, because
that might violate atomicity.}

\textbf{False.}  A thread can, for example, issue a blocking I/O operation
while retaining a mutex.  Other threads then remain excluded, so the mutual
exclusion, and hence the critical section's atomicity with respect to those
threads, is not violated.  Doing this is often poor design because it can
cause long delays, lock convoys, or deadlock.  Some special locks, such as
kernel spinlocks, impose a separate rule that their holders must not sleep,
but that is not a property of every lock.

\end{enumerate}

\subsection*{Major 2020 --- Q5: Ticket lock}

\emph{Ten threads correctly use a ticket lock whose \texttt{lock} operation
atomically fetches and increments \texttt{ticket}, waits for
\(\texttt{turn}=\texttt{myturn}\), and whose \texttt{unlock} increments
\texttt{turn}.  Determine whether each given
\((\texttt{ticket},\texttt{turn})\) pair is possible.}

Ignoring integer wraparound, correct use maintains
\[
0\leq \texttt{ticket}-\texttt{turn}\leq 10.
\]
The difference is exactly the number of threads that have taken a ticket but
have not yet completed the corresponding unlock: one may be in the critical
section and the rest are waiting.

\begin{enumerate}
\item[(a)] \((0,0)\): \textbf{possible}.  This is the initialized state,
before any acquisition.

\item[(b)] \((0,1)\): \textbf{impossible}.  It would mean that an unlock had
advanced \texttt{turn} although no ticket had ever been issued.

\item[(c)] \((1,0)\): \textbf{possible}.  One thread has received ticket 0
and is either about to enter or is in the critical section; it has not yet
unlocked.

\item[(d)] \((150,100)\): \textbf{impossible}.  There would be
\(150-100=50\) outstanding acquisitions, but only ten correctly using
threads can each have at most one outstanding acquisition.

\item[(e)] \((1005,1000)\): \textbf{possible}.  One thousand acquisitions
have completed and five of the ten threads currently have outstanding
tickets.
\end{enumerate}

\subsection*{Major 2020 --- Q6: Racy linked-list insertion}

\emph{Three threads insert keys 1, 2, and 3 into an initially empty list,
without synchronization, using}

\begin{verbatim}
n->next = head;
head = n;
\end{verbatim}

\emph{where allocation always succeeds and is thread-safe.  Give every
possible reachable list length and shape.}

Under the intended interleaving model, pointer reads and writes are atomic.
Two or more threads can read the same old value of \texttt{head} in the
first assignment.  Their later writes to \texttt{head} then overwrite one
another, making some allocated nodes unreachable.

\begin{itemize}
\item Length 1:
\[
1,\qquad 2,\qquad 3.
\]
For example, all three threads can read \texttt{NULL} before any of them
writes \texttt{head}; whichever writes \texttt{head} last is the only
reachable node.

\item Length 2 (every ordered pair of distinct keys):
\[
\begin{array}{lll}
1\mathbin{\to}2,&1\mathbin{\to}3,&
2\mathbin{\to}1,\\
2\mathbin{\to}3,&3\mathbin{\to}1,&
3\mathbin{\to}2.
\end{array}
\]
The remaining node was allocated but lost.  For example, threads 2 and 3
may both prepare a node pointing to \texttt{NULL}; thread 2 publishes its
node, thread 1 records node 2 as its successor, thread 3 temporarily
overwrites the head, and thread 1 finally publishes node 1.  The reachable
list is \(1\to2\).

\item Length 3 (every permutation):
\[
\begin{array}{lll}
1\mathbin{\to}2\mathbin{\to}3,&
1\mathbin{\to}3\mathbin{\to}2,&
2\mathbin{\to}1\mathbin{\to}3,\\
2\mathbin{\to}3\mathbin{\to}1,&
3\mathbin{\to}1\mathbin{\to}2,&
3\mathbin{\to}2\mathbin{\to}1.
\end{array}
\]
Each is obtained by a suitable serial ordering (or an equivalent
interleaving) of the three insertions.
\end{itemize}

Length 0 is impossible because at least one of the three writes to
\texttt{head} occurs.  No cycle can be formed: an edge from node \(i\) to
node \(j\) requires \(j\) to have been published before \(i\)'s first
assignment; the resulting strict ``published earlier'' order cannot form a
cycle.  Strictly speaking, unsynchronized conflicting accesses to
\texttt{head} are a C data race and therefore have undefined behaviour;
the lists above are the complete outcomes in the sequentially consistent,
atomic-pointer model intended by the question.

\subsection*{Major 2020 --- Q7: Synchronizing a one-doctor clinic}

\emph{There is one doctor, an infinite waiting room, and one shared buffer.
A patient calls \texttt{consultDoctor()} to write consultation information
and later calls \texttt{noteTreatment()} to read the treatment.  The doctor
repeatedly calls \texttt{treatPatient()}, which reads the consultation and
writes the treatment.  Use only semaphores so patients do not overwrite one
another, the doctor reads only new valid information, and a patient reads
only a completed treatment.}

Three semaphores suffice:
\[
\begin{array}{c|c|l}
\mbox{name} & \mbox{initial value} & \mbox{purpose}\\ \hline
\texttt{office} & 1 & \mbox{only one patient may use the shared buffer}\\
\texttt{patientInfo} & 0 & \mbox{a complete, new consultation is available}\\
\texttt{treatmentReady} & 0 & \mbox{the doctor has written the treatment}
\end{array}
\]
Here \(P\) is semaphore wait and \(V\) is semaphore signal.

\begin{verbatim}
Patient()
{
    P(office);             /* wait outside; reserve the buffer */
    consultDoctor();       /* write this patient's information */
    V(patientInfo);        /* publish it to the doctor */

    P(treatmentReady);     /* wait until doctor writes result */
    noteTreatment();       /* read this patient's treatment */
    V(office);             /* next patient may use the buffer */
}

Doctor()
{
    while (1) {
        P(patientInfo);    /* never read an empty/stale request */
        treatPatient();    /* read request and write treatment */
        V(treatmentReady);
    }
}
\end{verbatim}

\texttt{office} is held by a patient for the complete request--response
exchange, so another patient cannot overwrite either phase of the shared
buffer.  The two zero-valued semaphores provide the request and response
handshakes, including the required memory-ordering points.  The semaphore's
wait queue represents the unbounded waiting room.

\subsection*{Minor 2 2016 --- Q1: Bounded buffer using \(n-1\) vs.\ \(n\) slots}
\textbf{Question.}
First, give a solution (algorithm) which allows at most \(n-1\) entries in
the buffer, of the producer--consumer problem with a finite buffer.
Secondly, modify the algorithm to remedy this deficiency.

\textbf{Solution.}
Let \texttt{buffer[0..n-1]} be a circular buffer.  Initially,
\(\texttt{in}=\texttt{out}=0\).  For one producer and one consumer, the
following algorithm reserves one array position in order to distinguish
``empty'' from ``full'':

\begin{verbatim}
/* Producer */
for (;;) {
    item = produce_item();
    while ((in + 1) % n == out)
        ;                       /* full: spin */
    buffer[in] = item;
    in = (in + 1) % n;
}

/* Consumer */
for (;;) {
    while (in == out)
        ;                       /* empty: spin */
    item = buffer[out];
    out = (out + 1) % n;
    consume_item(item);
}
\end{verbatim}

Here \(\texttt{in}=\texttt{out}\) denotes an empty buffer, whereas
\((\texttt{in}+1)\bmod n=\texttt{out}\) denotes a full buffer.  Thus at
most \(n-1\) items can be stored: one physical slot is deliberately
wasted.  (With multiple producers or multiple consumers, the index
updates would additionally need mutual exclusion.)

To use all \(n\) slots, keep empty and full states separately using
counting semaphores.  Let
\[
  \texttt{mutex}=1,\qquad \texttt{empty}=n,\qquad \texttt{full}=0,
  \qquad \texttt{in}=\texttt{out}=0.
\]
The corrected bounded-buffer algorithm is:

\begin{verbatim}
/* Producer */
for (;;) {
    item = produce_item();
    sem_wait(&empty);
    sem_wait(&mutex);
    buffer[in] = item;
    in = (in + 1) % n;
    sem_post(&mutex);
    sem_post(&full);
}

/* Consumer */
for (;;) {
    sem_wait(&full);
    sem_wait(&mutex);
    item = buffer[out];
    out = (out + 1) % n;
    sem_post(&mutex);
    sem_post(&empty);
    consume_item(item);
}
\end{verbatim}

\texttt{empty} now records the number of free slots and \texttt{full}
records the number of occupied slots, so the state with
\(\texttt{full}=n\) is unambiguous.  Hence all \(n\) slots can be used.
\texttt{mutex} protects the buffer and indices, so this version also
works with multiple producers and consumers.


\clearpage
% ===========================================================================
\section{Lecture 15: Beyond physical memory (swapping)}
\renewcommand{\lecno}{15}
% ===========================================================================

\keepspace{17\baselineskip}
\begin{Qn}{T/F}{Swapping basics}
State true or false, with a one-line justification.
\begin{enumerate}[(a)]
  \item With swapping, a process can run even if its virtual address space
        is larger than physical memory.
  \item A page fault means the process accessed an illegal address.
  \item The hardware clears the use (reference) bits periodically.
  \item For any reference string, LRU with 4 frames never has more faults
        than LRU with 3 frames.
  \item Page faults are handled entirely by hardware on machines with a
        hardware-managed TLB.
  \item Evicting a clean page requires no disk write.
  \item Prepaging always reduces the number of page faults.
\end{enumerate}
\end{Qn}
\begin{Sol}
\begin{enumerate}[(a)]
  \item \textbf{True.} Only the pages in use need to be in memory; the rest
        live on disk (present=0). That is the whole illusion.
  \item \textbf{False.} Page fault = valid=1 but present=0: a \emph{legal}
        page that is on disk. An illegal address (valid=0) gives a
        segmentation fault.
  \item \textbf{False.} Hardware only \emph{sets} the use bit on a
        reference; the OS clears it (e.g.\ the clock hand).
  \item \textbf{True.} LRU has the stack property: the 3-frame contents are
        always a subset of the 4-frame contents, so every 3-frame hit is a
        4-frame hit. (FIFO lacks this: Belady's anomaly.)
  \item \textbf{False.} Page faults always trap into the OS: disk I/O is
        too slow and complex for hardware.
  \item \textbf{True.} Its copy on disk is still correct, so it is simply
        discarded; only dirty pages are written back.
  \item \textbf{False.} Wrong predictions waste disk I/O and the prefetched
        pages can evict pages that were about to be used, \emph{causing}
        extra faults.
\end{enumerate}
\end{Sol}

\begin{Qn}{Concept}{Valid vs.\ present}
A process's page table is:
\begin{center}
\renewcommand{\arraystretch}{1.08}
\begin{tabular}{@{}ccccc@{}}
  \toprule
  \textbf{VPN} & \textbf{PFN / disk block} & \textbf{valid} & \textbf{prot} & \textbf{present} \\
  \midrule
  0 & 7  & 1 & \texttt{r-x} & 1 \\
  1 & -- & 0 & --           & -- \\
  2 & 41 & 1 & \texttt{rw-} & 0 \\
  3 & 3  & 1 & \texttt{rw-} & 1 \\
  \bottomrule
\end{tabular}
\end{center}
For each access, say what happens and what the OS does: (a) load from
VPN 0, (b) store to VPN 2, (c) load from VPN 1, (d) store to VPN 0.
\end{Qn}
\begin{Sol}
\begin{enumerate}[(a)]
  \item Present, readable: normal translation to PFN 7. No OS involvement
        (after a TLB miss at most).
  \item valid=1, present=0: \textbf{page fault}. The 41 is a \emph{disk
        block}. The OS picks a free frame (evicting a victim if none, and
        writing it back first if dirty), reads block 41 into it, sets
        PFN=new frame and present=1, and restarts the store. Since it is a
        store, the page becomes dirty.
  \item valid=0: not part of the address space $\Rightarrow$
        \textbf{segmentation fault}; the process is likely killed.
  \item Present but \texttt{r-x}: a \textbf{protection fault} (write to a
        read-only code page); the OS treats it as an illegal access.
\end{enumerate}
\end{Sol}

\begin{Qn}{Concept}{Order the page-fault steps}
Put these in order and say whether hardware (HW) or OS does each:
(i) update PTE: PFN and present=1;
(ii) restart the faulting instruction;
(iii) look up the PTE for the virtual address;
(iv) read the page from disk into a frame;
(v) trap into the OS because present=0;
(vi) find the disk address in the PTE;
(vii) choose a victim frame and write it back if dirty.
While (iv) is in progress, what should the scheduler do?
\end{Qn}
\begin{Sol}
(iii) HW (TLB miss, page-table walk) $\to$ (v) HW raises the fault $\to$
(vi) OS $\to$ (vii) OS (only if no free frame) $\to$ (iv) OS issues
disk I/O $\to$ (i) OS $\to$ (ii) HW re-executes the instruction when the
process runs again.

During (iv) the disk read takes milliseconds: the faulting process is
\textbf{blocked} and the scheduler runs \textbf{another ready process}.
When the disk interrupt arrives, the process becomes ready again.
\end{Sol}

\begin{Qn}{Concept}{Restarting instructions}
Why does demand paging need \emph{precise interrupts}? Explain the problem
with an instruction like \texttt{move +(SP), R2} (auto-increment), and
state what the hardware must do.
\end{Qn}
\begin{Sol}
A page fault can occur \emph{in the middle of} an instruction (on fetch or
on the data access), and the fault must be transparent: after the OS loads
the page, the instruction must run again as if nothing happened. Precise
interrupts stop the pipeline so that all earlier instructions have
completed and the faulting one (and later ones) can be restarted from
scratch.

\texttt{move +(SP), R2} increments \texttt{SP} and loads from memory. If
the increment happens and then the load faults, re-running the whole
instruction increments \texttt{SP} \emph{twice}. The hardware must
\textbf{track side effects} and undo them before the restart (or be able
to resume the instruction where it stopped).
\end{Sol}

\begin{Qn}{Concept}{Locality and policies}
(a) Name the two kinds of locality and give a code example of each.
(b) Why does it make sense to spend thousands of instructions choosing a
victim page? (c) Give one access pattern where prepaging helps and one
where it hurts.
\end{Qn}
\begin{Sol}
(a) \textbf{Temporal}: the same data is reused soon, e.g.\ a loop counter
\texttt{i} or \texttt{sum} accessed every iteration. \textbf{Spatial}:
nearby addresses are used soon, e.g.\ scanning \texttt{a[0], a[1], \ldots},
or sequential instructions.

(b) A page fault costs a disk access (\textbf{milliseconds} = millions of
instructions). Even a few thousand instructions of policy code are
negligible if they avoid a single extra fault. The goal is to minimize
the number of faults.

(c) Helps: \textbf{sequential} access (reading a file or array page by
page: fetch \(P+1\) with \(P\)). Hurts: random access (the predicted pages
are never used; they waste I/O and evict useful pages).
\end{Sol}

\keepspace{24\baselineskip}
\begin{Qn}{Trace}{OPT, FIFO, LRU}
Reference string \texttt{7 0 1 2 0 3 0 4 2 3 0 3 2}, \textbf{3 frames},
initially empty. For each of OPT, FIFO and LRU give the frame contents
after every reference, mark misses, and give the total number of misses.
\end{Qn}
\begin{Sol}
\begin{center}
\small
\renewcommand{\arraystretch}{1.02}
\setlength{\tabcolsep}{4pt}
\begin{tabular}{@{}c|lc|lc|lc@{}}
  \toprule
  & \multicolumn{2}{c|}{\textbf{OPT}} & \multicolumn{2}{c|}{\textbf{FIFO} (oldest first)} & \multicolumn{2}{c}{\textbf{LRU} (least recent first)} \\
  \textbf{Ref} & frames & miss & queue & miss & order & miss \\
  \midrule
  7 & 7       & M & 7       & M & 7       & M \\
  0 & 0 7     & M & 7 0     & M & 7 0     & M \\
  1 & 0 1 7   & M & 7 0 1   & M & 7 0 1   & M \\
  2 & 0 1 2 (evict 7) & M & 0 1 2 & M & 0 1 2 & M \\
  0 & 0 1 2   &   & 0 1 2   &   & 1 2 0   &   \\
  3 & 0 2 3 (evict 1) & M & 1 2 3 & M & 2 0 3 & M \\
  0 & 0 2 3   &   & 2 3 0   & M & 2 3 0   &   \\
  4 & 2 3 4 (evict 0) & M & 3 0 4 & M & 3 0 4 & M \\
  2 & 2 3 4   &   & 0 4 2   & M & 0 4 2   & M \\
  3 & 2 3 4   &   & 4 2 3   & M & 4 2 3   & M \\
  0 & 0 2 3 (evict 4) & M & 2 3 0 & M & 2 3 0 & M \\
  3 & 0 2 3   &   & 2 3 0   &   & 2 0 3   &   \\
  2 & 0 2 3   &   & 2 3 0   &   & 0 3 2   &   \\
  \midrule
  & \multicolumn{2}{r|}{\textbf{7}} & \multicolumn{2}{r|}{\textbf{10}} & \multicolumn{2}{r}{\textbf{9}} \\
  \bottomrule
\end{tabular}
\end{center}
OPT evicts the page used furthest in the future: at ``2'' both 7 and 1
are never used again (a tie: evicting either is optimal and gives the same
total), and we evict 7; at ``3'' it evicts 1 (never used again); at ``4''
it evicts 0 (next use is furthest); at the second ``0'' it evicts 4.
OPT \(\le\) LRU \(\le\) FIFO here, as expected (not guaranteed in
general for LRU vs.\ FIFO).
\end{Sol}

\begin{Qn}{Trace}{Belady's anomaly, and why LRU is immune}
Reference string \texttt{1 2 3 4 1 2 5 1 2 3 4 5}.
(a) Count FIFO misses with 3 and with 4 frames.
(b) Count LRU misses with 3 and with 4 frames.
(c) Explain the difference using the stack property.
\end{Qn}
\begin{Sol}
(a) FIFO: \textbf{9} misses with 3 frames, \textbf{10} with 4 frames:
more memory, more faults (Belady's anomaly). (This is the lecture's
\texttt{A B C D A B E A B C D E} example with letters renamed.)

(b) LRU: \textbf{10} misses with 3 frames, \textbf{8} with 4 frames.

(c) LRU with \(k\) frames always holds the \(k\) most recently used pages;
with \(k+1\) frames it holds the \(k+1\) most recently used, a superset. So
every hit with \(k\) frames is a hit with \(k+1\) frames, and faults can
only go down. FIFO's contents depend on arrival order, and the 3-frame set
need not be a subset of the 4-frame set (e.g.\ after ``5'': FIFO-3 =
\(\{1,2,5\}\), FIFO-4 = \(\{2,3,4,5\}\)), so the anomaly can occur.
\end{Sol}

\keepspace{26\baselineskip}
\begin{Qn}{Trace}{Clock algorithm}
4 frames (0--3), initially empty, clock hand at frame 0. On a fault, an
empty frame is used first (use bit set to 1); once all frames are full,
the clock algorithm picks the victim, and the new page's use bit is set to
1 and the hand advances past it. A hit sets the page's use bit to 1.
Reference string: \texttt{1 2 3 4 1 5 2 6 1 3}.
Give the victim on every eviction and the total number of faults. What
happens at the reference to 5, and why?
\end{Qn}
\begin{Sol}
\begin{center}
\small
\renewcommand{\arraystretch}{1.05}
\begin{tabular}{@{}clll@{}}
  \toprule
  \textbf{Ref} & \textbf{Result} & \textbf{Frames 0--3} & \textbf{Use bits, hand after} \\
  \midrule
  1--4 & 4 misses (fill) & 1 2 3 4 & 1 1 1 1, hand 0 \\
  1 & hit & 1 2 3 4 & 1 1 1 1, hand 0 \\
  5 & miss, \textbf{evict 1} & 5 2 3 4 & 1 0 0 0, hand 1 \\
  2 & hit & 5 2 3 4 & 1 1 0 0, hand 1 \\
  6 & miss, \textbf{evict 3} & 5 2 6 4 & 1 0 1 0, hand 3 \\
  1 & miss, \textbf{evict 4} & 5 2 6 1 & 1 0 1 1, hand 0 \\
  3 & miss, \textbf{evict 2} & 5 3 6 1 & 0 1 1 1, hand 2 \\
  \bottomrule
\end{tabular}
\end{center}
Total: \textbf{8 faults} (2 hits). At ``5'' all use bits are 1, so the
hand clears all four in one full sweep and evicts the frame it started at
(frame 0, page 1), even though 1 was just used. When every bit is set,
clock degenerates to \textbf{FIFO}. At ``6'' the hand clears frame 1 (page
2 was re-used) and evicts page 3, which was not used since the sweep: the
LRU-like behaviour. (For comparison: OPT 6, FIFO 8, LRU 8 faults.)
\end{Sol}

\begin{Qn}{Concept}{Enhanced clock with dirty bits}
Frames hold pages with (use, dirty) bits: P = (1,1), Q = (0,1), R = (1,0),
S = (0,1), T = (1,0), in clock order starting at the hand. (a) Which page
is evicted using the class rule? (b) Why is (0,1) preferred over (1,0),
even though evicting (0,1) costs a disk write?
\end{Qn}
\begin{Sol}
(a) Classes from best to worst: (0,0), (0,1), (1,0), (1,1). No page is
(0,0); the lowest non-empty class is (0,1), containing Q and S. The
\textbf{first} in clock order is \textbf{Q}.

(b) ``Recently used'' matters more than ``dirty''. A (1,0) page was used
recently and will probably be needed again, so evicting it likely costs a
whole page \emph{fault} later (disk read + process blocked). A (0,1) page
costs one write now, but is probably not needed soon.
\end{Sol}

\begin{Qn}{Trace}{Nth chance}
Nth chance with \(N=2\): each sweep, a used page gets
\texttt{notInUseSince = 0}; an unused page with \texttt{notInUseSince} \(<
N\) gets it incremented; otherwise it is reclaimed. Page X is used before
sweep 1 and never again. During which sweep is X reclaimed? How does
\(N\) trade accuracy against cost, and what is \(N=1\)?
\end{Qn}
\begin{Sol}
Sweep 1: X used $\Rightarrow$ 0. Sweep 2: unused, \(0<2\) $\Rightarrow$ 1.
Sweep 3: \(1<2\) $\Rightarrow$ 2. Sweep 4: \(2 \not< 2\) $\Rightarrow$
\textbf{reclaimed in sweep 4} (after \(N=2\) full sweeps without use).

Large \(N\): finer distinction between old and very old pages, closer to
true LRU, but more sweeps (more work) before a victim is found. Small \(N\)
is cheaper. The ordinary \textbf{clock} reclaims a page the first time a
sweep finds it unused; with this exact pseudocode that is \(N = 0\) (for
\(N = 1\), X would be reclaimed in sweep 3).
\end{Sol}

\begin{Qn}{Concept}{Implementing perfect LRU}
Compare software LRU (ordered list) and hardware LRU (timestamp per page):
what is slow in each? For a 4\,GB machine with 4\,KB pages, how many
timestamps must be scanned to find a victim? What do real systems do?
\end{Qn}
\begin{Sol}
Software list: move the page to the front on \emph{every memory
reference}: \textbf{slow references}, fast replacement (take the back).
Hardware timestamps: storing the clock on a reference is fast, but finding
the victim means \textbf{scanning all} pages: \(4\,\text{GB} / 4\,\text{KB} =
2^{32}/2^{12} = 2^{20} \approx 10^6\) registers. Real systems do not
implement perfect LRU; they approximate it (clock with use bits, Nth
chance), since LRU itself is only an approximation of OPT: find an
\emph{old} page, not necessarily the oldest.
\end{Sol}

% ===========================================================================
\section{Lecture 16: Concurrency and threads}
\renewcommand{\lecno}{16}
% ===========================================================================

\begin{Qn}{Concept}{Shared or private?}
For two threads of the same process, classify each as \textbf{shared} or
\textbf{private (per-thread)}: (a) a global variable, (b) a local variable
of the thread function, (c) a \texttt{malloc}'d object whose pointer is
stored in a global, (d) the PC/instruction pointer, (e) the stack pointer,
(f) the page table / PTBR value, (g) open file descriptors, (h) the PID,
(i) \texttt{\%eax}.
\end{Qn}
\begin{Sol}
Shared: (a), (c), (f) (same address space), (g), (h).
Private: (b) (on the thread's own stack), (d), (e), (i) (each thread has
its own register set, saved/restored on a switch).
Note: (b) is private in use but still in the shared address space; another
thread could reach it through a pointer.
\end{Sol}

\begin{Qn}{T/F}{Processes vs.\ threads}
\begin{enumerate}[(a)]
  \item Switching between two threads of the same process changes the PTBR.
  \item Threads communicate through \texttt{pipe()} because they cannot
        share memory.
  \item \texttt{pthread\_join} is the thread analogue of \texttt{wait()}.
  \item Making \texttt{counter} \texttt{volatile} makes
        \texttt{counter = counter + 1} atomic.
  \item If a multi-threaded program passes all tests, it has no race
        conditions.
\end{enumerate}
\end{Qn}
\begin{Sol}
(a) \textbf{False}: same address space, same page table; that is why the
switch is cheaper (no TLB/cache disruption from a new address space).
(b) \textbf{False}: threads share the address space; communication is
plain memory reads/writes (processes needed pipes and copying).
(c) \textbf{True}.
(d) \textbf{False}: \texttt{volatile} only forces real loads/stores to
memory; the update is still load, add, store.
(e) \textbf{False}: the bad interleaving may simply not have happened in
the test runs. Assume a malicious scheduler.
\end{Sol}

\begin{Qn}{Predict}{Thread output orders}
\texttt{main} prints \texttt{main: begin}, creates \(p_1\) (prints
\texttt{A}) and \(p_2\) (prints \texttt{B}), then calls
\texttt{pthread\_join(p2)}, \texttt{pthread\_join(p1)}, and prints
\texttt{main: end}. List all possible outputs. Does swapping the join order
change anything?
\end{Qn}
\begin{Sol}
\texttt{main: begin}, then \texttt{A} and \texttt{B} in \textbf{either
order}, then \texttt{main: end}: exactly 2 outputs. \texttt{main: end}
always comes last because both joins return only after both threads finish.
Join order does not affect what the threads print or when; it only decides
which finished thread main notices first.
\end{Sol}

\begin{Qn}{Bad schedule}{Returning a value from a thread}
What is wrong with this thread function, and how did the lecture's version
avoid it?
\begin{lstlisting}[numbers=none]
void *mythread(void *arg) {
    myret_t r;  r.x = 1;  r.y = 2;
    return (void *) &r;
}
\end{lstlisting}
\end{Qn}
\begin{Sol}
\texttt{r} is a local variable on the \textbf{thread's own stack}. When the
thread returns, that stack frame is gone, so \texttt{pthread\_join}
receives a dangling pointer. The lecture allocated the return value with
\texttt{Malloc} (heap), which outlives the thread; \texttt{main} then frees
it after reading it.
\end{Sol}

\begin{Qn}{Trace}{Register and memory state}
\texttt{counter} at \texttt{0x9cd4} starts at 50. Each thread runs
\texttt{mov 0x9cd4,\%eax}; \texttt{add \$0x1,\%eax};
\texttt{mov \%eax,0x9cd4}. Schedule: T1 runs 2 instructions, switch, T2
runs all 3, switch, T1 runs its last instruction. Give \texttt{\%eax} of
each thread and the memory value after every step, and the final value.
\end{Qn}
\begin{Sol}
\begin{center}
\small
\begin{tabular}{@{}llccc@{}}
  \toprule
  \textbf{Runs} & \textbf{Instruction} & \textbf{T1 \%eax} & \textbf{T2 \%eax} & \textbf{0x9cd4} \\
  \midrule
  T1 & \texttt{mov 0x9cd4,\%eax} & 50 & -- & 50 \\
  T1 & \texttt{add} & 51 & -- & 50 \\
  \multicolumn{5}{c}{\textit{switch: T1's \%eax = 51 saved in its control block}} \\
  T2 & \texttt{mov 0x9cd4,\%eax} & (51) & 50 & 50 \\
  T2 & \texttt{add} & (51) & 51 & 50 \\
  T2 & \texttt{mov \%eax,0x9cd4} & (51) & 51 & 51 \\
  \multicolumn{5}{c}{\textit{switch: T1 restored with \%eax = 51}} \\
  T1 & \texttt{mov \%eax,0x9cd4} & 51 & (51) & \textbf{51} \\
  \bottomrule
\end{tabular}
\end{center}
Final \textbf{51} instead of 52: one increment is lost because both threads
loaded 50 before either stored.
\end{Sol}

\begin{Qn}{Predict}{How many interleavings, which results?}
Thread 1 adds 1 and Thread 2 adds 2 to shared \texttt{x} (each: load, add,
store; \texttt{x} = 0 initially). (a) How many distinct interleavings of
the 6 instructions are there? (b) Which final values are possible, and in
how many interleavings each? (c) In general, how many interleavings do
two threads with \(m\) and \(n\) atomic steps have?
\end{Qn}
\begin{Sol}
(a) Choose the 3 positions of Thread 1's instructions among 6:
\(\binom{6}{3} = \mathbf{20}\).

(b) Final value \textbf{3} (correct) in only \textbf{2} interleavings: the
two serial ones (one thread entirely before the other's load). Otherwise
both threads load the old value 0 and the \textbf{last store wins}: final
\textbf{2} (T2 stores last) in \textbf{9}, final \textbf{1} (T1 stores
last) in \textbf{9}. So \(2+9+9 = 20\).

(c) \(\binom{m+n}{m}\) (each thread's own order is fixed). E.g.\ steps
\(A;B;C\) and \(D;E\) give \(\binom{5}{2} = 10\) traces (a past exam
question).
\end{Sol}

\begin{Qn}{Predict}{Range of a racy counter (hard)}
Two threads each execute \texttt{counter = counter + 1} twice (load, add,
store each time); \texttt{counter} starts at 0. What final values are
possible? Give a schedule achieving the minimum.
\end{Qn}
\begin{Sol}
Possible: \textbf{2, 3, 4}. (Maximum 4: no overlap.)

Minimum 2 schedule: T1 loads 0. T2 does a full increment (counter = 1),
then loads 1 for its second increment. T1 adds and stores 1, then does its
whole second increment (load 1, store 2). Finally T2 adds 1 to its stale
register value 1 and stores \textbf{2}. Three of four increments are
``partly lost'', but the last increment of each thread still counts once,
which is why the result cannot go below 2. (With \(N \ge 2\) iterations each, the
final value can be anything from 2 to \(2N\); with \(N = 1\) it is 1 or 2.)
\end{Sol}

\begin{Qn}{Concept}{Definitions}
Define: race condition, critical section, mutual exclusion, atomicity,
non-determinism. Then explain why ``imagine the scheduler is malicious'' is
the right way to reason about concurrent code.
\end{Qn}
\begin{Sol}
\textbf{Race condition}: the result depends on the timing (interleaving) of
threads. \textbf{Critical section}: code that accesses shared data and must
not be executed by more than one thread at a time. \textbf{Mutual
exclusion}: if one thread is in the critical section, others cannot be
(unrelated work is fine). \textbf{Atomicity}: a group of instructions
executes as an uninterruptable unit (all or nothing, no interleaving
inside). \textbf{Non-determinism}: same inputs, different results on
different runs.

A bug shows up only under some schedules; tests may never hit them. If code
is correct even when the scheduler picks the worst possible switch points,
it is correct for every schedule.
\end{Sol}

\begin{Qn}{Concept}{Why threads?}
(a) Why can't a single-threaded process use a second core? (b) Give two
costs of building an application from many communicating processes instead
of threads. (c) Why does CPU frequency stagnation make this matter?
\end{Qn}
\begin{Sol}
(a) A process with one thread has one stream of execution; it can run on
only one CPU at a time, and cannot overlap its own I/O with computation.
(b) Cumbersome programming (separate address spaces, explicit messages),
copying overheads (data goes through the kernel via \texttt{pipe()}),
expensive context switches (new address space: PTBR change, TLB and cache
disruption). (c) Clock frequency flattened around 2005; new chips add
\textbf{more cores} rather than faster ones, so programs only get faster
through concurrent execution.
\end{Sol}

% ===========================================================================
\section{Lecture 17: Locks}
\renewcommand{\lecno}{17}
% ===========================================================================

\begin{Qn}{Predict}{Racy linked-list insert}
The list is \texttt{head} \(\to\) X. Threads A and B each run
\texttt{new->next = L->head; L->head = new;} (no lock) to insert nodes a
and b. List every possible final list, and which node (if any) is lost.
\end{Qn}
\begin{Sol}
\begin{itemize}
  \item \(a \to b \to X\): B completes, then A runs (A reads head = b).
  \item \(b \to a \to X\): A completes, then B runs.
  \item \(a \to X\), \textbf{b lost}: both read head = X, B writes head =
        b, then A writes head = a.
  \item \(b \to X\), \textbf{a lost}: both read X, A writes first, B last.
\end{itemize}
In the lost cases the orphan still points to X but nothing points to it.
The fix is a lock around the two lines (Approach \#2: the smallest
critical section).
\end{Sol}

\begin{Qn}{Concept}{Where should the lock go?}
For \texttt{List\_Insert} (malloc, set key, set next, set head) and
\texttt{List\_Lookup} (walk the list):
(a) which lines of Insert must be inside the critical section, and why
can \texttt{malloc} be outside?
(b) Approach \#1 locks the whole of Lookup but unlocks only before
\texttt{return 0}. What bug is that?
(c) Why may Lookup skip the lock entirely if there is no
\texttt{List\_Delete()}?
\end{Qn}
\begin{Sol}
(a) Only \texttt{new->next = L->head; L->head = new;} (the read and write
of shared \texttt{head}). \texttt{malloc} and \texttt{new->key = key} touch
a node no other thread can see yet (assuming \texttt{malloc} is
thread-safe). Bonus: a failed \texttt{malloc} then holds no lock.

(b) The \texttt{return 1} path (key found) returns while still
\textbf{holding the lock}: every later Insert/Lookup blocks forever. Every
return path must unlock.

(c) Lookup only reads; Insert publishes a fully built node with a single
store to \texttt{head}, so a reader sees either the old or new head, and
nodes are never freed under it. With deletes, a lookup could follow a
pointer into a freed node.
\end{Sol}

\begin{Qn}{T/F}{Disabling interrupts}
\begin{enumerate}[(a)]
  \item Disabling interrupts gives mutual exclusion on a multiprocessor.
  \item A user program that disables interrupts can monopolize the CPU.
  \item The OS kernel never disables interrupts.
\end{enumerate}
\end{Qn}
\begin{Sol}
(a) \textbf{False}: it only stops preemption on the \emph{current} CPU;
threads on other CPUs can still enter the critical section.
(b) \textbf{True}: no timer interrupt means the OS never regains control
(and it also misses I/O interrupts). That is why it must be privileged.
(c) \textbf{False}: the OS itself briefly masks interrupts around its own
data structures; inside the kernel the trust problem disappears.
\end{Sol}

\begin{Qn}{Bad schedule}{The single-flag lock}
\texttt{acquire: while (*lock) ; *lock = true;} --
\texttt{release: *lock = false;}. Give a 2-thread schedule that breaks
mutual exclusion, and name the exact property that is missing.
\end{Qn}
\begin{Sol}
\texttt{*lock} = 0. T1: \texttt{while(*lock)} sees 0, exits loop.
\emph{Switch.} T2: sees 0, exits, sets \texttt{*lock = 1}, enters. \emph{Switch.}
T1: sets \texttt{*lock = 1}, enters. \textbf{Both are in the critical
section.} Testing the lock and setting it are \textbf{not atomic}; a switch
between the test and the set breaks it.
\end{Sol}

\keepspace{20\baselineskip}
\begin{Qn}{Bad schedule}{Broken variants of Peterson's algorithm}
Correct Peterson (\texttt{tid} \(\in \{0,1\}\)):
\texttt{lock[tid]=true; turn=1-tid; while(lock[1-tid] \&\& turn==1-tid);}.
For each variant, say which property fails and give a schedule.
\begin{enumerate}[(a)]
  \item The two stores are swapped: \texttt{turn=1-tid; lock[tid]=true;}.
  \item No \texttt{turn}: \texttt{lock[tid]=true; while(lock[1-tid]);}.
  \item No \texttt{lock[]}: \texttt{turn=1-tid; while(turn==1-tid);}.
\end{enumerate}
\end{Qn}
\begin{Sol}
(a) \textbf{Mutual exclusion fails.} T0: \texttt{turn=1}. T1:
\texttt{turn=0}, \texttt{lock[1]=true}, tests \texttt{lock[0] \&\&
turn==0}: \texttt{lock[0]} is still false $\Rightarrow$ T1 enters. T0:
\texttt{lock[0]=true}, tests \texttt{lock[1] \&\& turn==1}: turn is 0
$\Rightarrow$ T0 enters. Both inside. (Announcing interest \emph{after}
giving away the turn lets the other thread miss it.)

(b) \textbf{Progress fails (deadlock).} T0: \texttt{lock[0]=true}. T1:
\texttt{lock[1]=true}. Now both spin forever, each waiting for the other's
flag. (Mutual exclusion still holds.)

(c) \textbf{Progress fails.} If only T0 wants the lock: \texttt{turn=1},
then it waits while \texttt{turn==1} forever; it can only enter after T1
also tries (strict alternation). Mutual exclusion holds.
\end{Sol}

\begin{Qn}{Concept}{Peterson: why it works, and its limits}
(a) Argue mutual exclusion for Peterson's algorithm. (b) Show bounded
waiting. (c) Why is it not used on modern CPUs?
\end{Qn}
\begin{Sol}
(a) Suppose both are inside. Both flags are true. Consider whichever thread
wrote \texttt{turn} \emph{last}, say T0 (it wrote \texttt{turn = 1}). T1's
flag was already true at that time and stays true while T1 is inside, so
T0's test \texttt{lock[1] \&\& turn==1} stays true and T0 cannot have
entered. Contradiction. (The last writer of \texttt{turn} waits.)

(b) If T1 is waiting and T0 leaves and immediately tries again, T0 sets
\texttt{turn = 1}, which makes T0 the one that waits. So T1 waits for at
most one critical section of T0.

(c) It assumes loads and stores happen in program order. Modern CPUs have
relaxed memory consistency (they may reorder them), so it can fail; real
locks use atomic hardware instructions instead.
\end{Sol}

\begin{Qn}{Fill}{Spinlocks from atomic instructions}
Fill in the blanks and explain why each works.
\begin{lstlisting}[numbers=none]
void init(lock_t *l)    { l->flag = __A__; }
void acquire(lock_t *l) { while (xchg(&l->flag, __B__) == __C__) ; }
void release(lock_t *l) { l->flag = __D__; }

void acquire2(lock_t *l) {   // with CompareAndSwap(addr, expected, new)
    while (CompareAndSwap(&l->flag, __E__, __F__) == __G__) ;
}
\end{lstlisting}
\end{Qn}
\begin{Sol}
A = 0, B = 1, C = 1, D = 0; E = 0, F = 1, G = 1.

\texttt{xchg} atomically writes 1 and returns the old value: old 0 means
the lock was free and we just took it; old 1 means it was held (writing 1
again changes nothing), so keep spinning. CAS sets the flag to 1
\emph{only if} it is 0 and returns the old value; we spin while it
returns 1. Both fix the flag lock because test and set happen in one atomic
instruction. (On a single CPU these spinlocks need a preemptive scheduler,
or the spinner never lets the holder run.)
\end{Sol}

\begin{Qn}{Concept}{Which lock gives what?}
Fill in yes/no for: Mutual exclusion (ME), Progress, Bounded waiting /
fairness, ``no wasted CPU while waiting'' --- for: disable interrupts (one
CPU), single flag, Peterson (2 threads), test-and-set spinlock, ticket
lock, ticket lock + \texttt{yield}, queue lock (\texttt{park}).
\end{Qn}
\begin{Sol}
\begin{center}
\small
\renewcommand{\arraystretch}{1.1}
\begin{tabular}{@{}lcccc@{}}
  \toprule
  \textbf{Lock} & \textbf{ME} & \textbf{Progress} & \textbf{Bounded/fair} & \textbf{No waste} \\
  \midrule
  Disable interrupts (1 CPU) & yes & yes & n/a & yes, but unsafe/privileged \\
  Single flag & \textbf{no} & yes & no & no (spins) \\
  Peterson (2 threads) & yes & yes & yes & no (spins) \\
  Test-and-set spinlock & yes & yes & \textbf{no} (can starve) & no \\
  Ticket lock & yes & yes & \textbf{yes} (FIFO order) & no \\
  Ticket + \texttt{yield} & yes & yes & yes & less (a switch per check) \\
  Queue lock + \texttt{park} & yes & yes & yes (queue order) & yes: waiters sleep \\
  \bottomrule
\end{tabular}
\end{center}
\end{Sol}

\begin{Qn}{Predict}{Possible ticket-lock states}
Five threads correctly use a ticket lock (\texttt{acquire}:
\texttt{myturn = FAA(\&ticket); while(turn != myturn);}; \texttt{release}:
\texttt{FAA(\&turn)}). Which (\texttt{ticket}, \texttt{turn}) pairs are
possible? (a) (0,0) (b) (3,3) (c) (2,3) (d) (9,3) (e) (8,3) (f) (12,7).
\end{Qn}
\begin{Sol}
Invariant: \(0 \le \texttt{ticket} - \texttt{turn} \le 5\). The difference
is the number of threads holding a ticket not yet released (one inside, the
rest spinning), and each thread has at most one outstanding ticket.
(a) possible (initial). (b) possible (3 completed, nobody waiting).
(c) \textbf{impossible} (turn \(>\) ticket: more releases than
acquisitions). (d) \textbf{impossible} (6 outstanding \(>\) 5 threads).
(e) possible (5 outstanding: 1 inside, 4 spinning). (f) possible (5
outstanding).
\end{Sol}

\begin{Qn}{Trace}{Ticket lock run}
\texttt{ticket = turn = 0}. Events: A lock, B lock, A unlock, C lock,
A lock, B unlock, C unlock, A unlock. After each event give
\texttt{ticket}, \texttt{turn}, each thread's \texttt{myturn}, and who is
in the critical section.
\end{Qn}
\begin{Sol}
\begin{center}
\small
\begin{tabular}{@{}lccl@{}}
  \toprule
  \textbf{Event} & \textbf{ticket} & \textbf{turn} & \textbf{Effect} \\
  \midrule
  A lock   & 1 & 0 & A: myturn 0 = turn $\Rightarrow$ \textbf{A in CS} \\
  B lock   & 2 & 0 & B: myturn 1, spins \\
  A unlock & 2 & 1 & \textbf{B in CS} \\
  C lock   & 3 & 1 & C: myturn 2, spins \\
  A lock   & 4 & 1 & A: myturn 3, spins (behind C) \\
  B unlock & 4 & 2 & \textbf{C in CS} \\
  C unlock & 4 & 3 & \textbf{A in CS} \\
  A unlock & 4 & 4 & nobody; lock free \\
  \bottomrule
\end{tabular}
\end{center}
Order of entry A, B, C, A: strictly by ticket. With a test-and-set lock A
could have re-grabbed the lock before C.
\end{Sol}

\begin{Qn}{Concept}{Spinning costs}
Four threads A--D share \textbf{one} CPU, time slice 20\,ms, context switch
0.5\,ms. A holds the lock and needs 5\,ms more of CPU in its critical
section; B, C, D want the lock. Round-robin order A, B, C, D, A, \ldots,
and A's slice expired just now. Estimate the time wasted before A releases
(a) with plain spinning, (b) with \texttt{yield()}. Generalize.
\end{Qn}
\begin{Sol}
(a) B, C, D each spin a full slice: \(3 \times 20 = 60\)\,ms wasted before A
runs again (then A releases after 5\,ms).
(b) Each of B, C, D runs only long enough to check and yield:
\(\approx 3 \times 0.5 = 1.5\)\,ms.

In general waste is \(O(\text{threads} \times \text{time\_slice})\) without
yield and \(O(\text{threads} \times \text{context\_switch})\) with yield. With
many threads even the yield version is costly (and the scheduler may keep
picking waiters), which motivates \textbf{blocking}: waiters leave the
ready queue entirely.
\end{Sol}

\begin{Qn}{Bad schedule}{Queue lock without \texttt{setpark}}
In the guard/queue lock, \texttt{acquire} does \texttt{qadd(q, tid);
guard = false; park();} when the lock is held. (a) Show the race with a
releasing thread. (b) How does \texttt{setpark()} fix it? (c) Why does
\texttt{release} not set \texttt{lock = false} when it unparks a waiter?
(d) Why is spinning on \texttt{guard} acceptable?
\end{Qn}
\begin{Sol}
(a) T1 (acquire): \texttt{qadd}, \texttt{guard = false}. \emph{Switch.}
T2 (release): takes guard, queue non-empty $\Rightarrow$
\texttt{unpark(T1)} (T1 is not asleep yet: the wakeup is lost),
\texttt{guard = false}. \emph{Switch.} T1: \texttt{park()} $\Rightarrow$
sleeps, possibly \textbf{forever}, although the lock was handed to it.

(b) T1 calls \texttt{setpark()} \emph{before} releasing guard. If an
\texttt{unpark} arrives between \texttt{setpark} and \texttt{park},
\texttt{park} returns immediately.

(c) The lock is \textbf{passed directly} to the woken thread: it returns
from \texttt{park} owning the lock without re-checking (it does not hold
guard there). If \texttt{lock} were set false, a third thread could grab it
first and both would be inside.

(d) \texttt{guard} is held only for the few instructions in
\texttt{acquire}/\texttt{release}, never during the user's critical
section, so the spin is very short.
\end{Sol}

\begin{Qn}{Concept}{Spin or block?}
A lock will be released after time \(t\); a context switch costs \(C\).
(a) How much CPU is wasted by spin-waiting? by blocking? (b) Decide for
\(t = 2\,\mu\)s, \(C = 10\,\mu\)s on a multiprocessor; and for \(t = 1\)\,ms.
(c) What should a waiter do on a \emph{uniprocessor}, whatever \(t\) is?
(d) What is the practical problem with this rule?
\end{Qn}
\begin{Sol}
(a) Spin: \(t\). Block: \(C\).
(b) \(t = 2\,\mu\text{s} < C\): \textbf{spin}. \(t = 1\,\text{ms} > C\):
\textbf{block}.
(c) \textbf{Always block / yield}: while the waiter runs, the holder cannot,
so spinning can never shorten the wait.
(d) It requires knowing \(t\) in advance (the future); predicting it costs
too much.
\end{Sol}

% ===========================================================================
\section{Lecture 18: Condition variables}
\renewcommand{\lecno}{18}
% ===========================================================================

\begin{Qn}{T/F}{CV semantics}
\begin{enumerate}[(a)]
  \item \texttt{pthread\_cond\_wait} may be called without holding the
        mutex.
  \item \texttt{wait} releases the mutex and sleeps atomically, and
        re-acquires the mutex before returning.
  \item A \texttt{signal} with no waiting thread is remembered and wakes the
        next thread that waits.
  \item After being signalled, the woken thread runs immediately and the
        condition is guaranteed to still hold.
  \item \texttt{broadcast} can always replace \texttt{signal} without
        breaking a correct \texttt{while}-based program.
\end{enumerate}
\end{Qn}
\begin{Sol}
(a) \textbf{False}: \texttt{wait} assumes the lock is held (it releases
it). (b) \textbf{True}. (c) \textbf{False}: a CV has no state besides its
queue; with no waiter the signal is lost (hence ``keep state''). (d)
\textbf{False}: Mesa semantics, a signal is only a hint; another thread
may run first and change the state (hence \texttt{while}). (e)
\textbf{True}: extra woken threads re-check and sleep again; only
performance suffers.
\end{Sol}

\begin{Qn}{Bad schedule}{Join attempts}
For each version, give a schedule (using the line labels) where the parent
waits forever, or argue it is correct.
\begin{enumerate}[(a)]
  \item Parent: \texttt{lock(x); wait(y); unlock(z)}.
        Child: \texttt{lock(a); signal(b); unlock(c)}.
  \item Parent: \texttt{lock(w); if (done==0)(x) wait(y); unlock(z)}.
        Child: \texttt{done=1 (a); signal (b)} (no lock).
  \item Parent as in (b). Child: \texttt{lock; signal; done=1; unlock}.
  \item Parent as in (b). Child: \texttt{lock; done=1; unlock; signal}.
\end{enumerate}
\end{Qn}
\begin{Sol}
(a) \textbf{Broken}: child a, b, c first (signal with nobody waiting is
lost), then parent x, y: \textbf{waits forever}. No state records that
the child finished.

(b) \textbf{Broken}: parent w, x (sees 0) $\to$ child a, b (signal lost:
parent not waiting yet) $\to$ parent y: sleeps forever. The child changes
the state outside the lock, in the gap between check and wait.

(c) \textbf{Correct.} Signal and \texttt{done = 1} are both inside the
lock. If the child runs first, the parent sees \texttt{done = 1}. If the
parent is waiting, it is woken, but can only return from \texttt{wait}
after re-acquiring the lock, i.e.\ after the child's unlock, by which time
\texttt{done = 1}. (Order inside a locked region does not matter to other
threads.)

(d) \textbf{Correct.} \texttt{done} is changed under the lock, so it
happens either before the parent's check (parent never waits) or after the
parent is already asleep (the later signal wakes it). Signalling without
the lock is not strictly needed here; it is still simplest to hold it.
\end{Sol}

\begin{Qn}{Concept}{Mesa semantics and \texttt{while}}
(a) What is Mesa semantics vs.\ Hoare semantics? (b) Why does Mesa force
\texttt{while} instead of \texttt{if}? (c) Give two separate reasons a
woken thread may find its condition false.
\end{Qn}
\begin{Sol}
(a) Mesa: \texttt{signal} just moves a waiter to the ready state; it runs
later, whenever scheduled. Hoare: the signalled thread runs
\emph{immediately}, with the state exactly as the signaller left it
(stronger, harder to build). Almost all real systems are Mesa.
(b) Between the signal and the woken thread re-acquiring the lock, other
threads can run and change the state; the woken thread must re-check.
(c) Another thread grabbed the lock first and consumed the state (e.g.\ a
second consumer took the item); \textbf{spurious wakeups} (a library may
wake a thread without a matching signal, or wake several).
\end{Sol}

\keepspace{30\baselineskip}
\begin{Qn}{Bad schedule}{Single CV with \texttt{if}: produce the trace}
Single buffer, one CV \texttt{cond}, producer \texttt{p1--p6} and consumer
\texttt{c1--c6} with \texttt{if (count == \ldots) wait}. With consumers
\(T_{c1}, T_{c2}\) and producer \(T_p\), give a trace ending with a consumer
calling \texttt{get()} on an empty buffer.
\end{Qn}
\begin{Sol}
\begin{center}
\small
\renewcommand{\arraystretch}{1.02}
\begin{tabular}{@{}lll@{}}
  \toprule
  \textbf{Step} & \textbf{count} & \textbf{What happens} \\
  \midrule
  \(T_{c1}\): c1, c2, c3 & 0 & nothing to get: \(T_{c1}\) sleeps \\
  \(T_p\): p1, p2, p4 & 1 & buffer now full \\
  \(T_p\): p5, p6 & 1 & signal: \(T_{c1}\) becomes \emph{ready} (not running) \\
  \(T_p\): p1, p2, p3 & 1 & buffer full: \(T_p\) sleeps \\
  \(T_{c2}\): c1, c2, c4 & 0 & \(T_{c2}\) sneaks in and grabs the data \\
  \(T_{c2}\): c5, c6 & 0 & signal wakes \(T_p\) \\
  \(T_{c1}\): c4 & 0 & \textbf{returns from wait, does not re-check: get() on empty!} \\
  \bottomrule
\end{tabular}
\end{center}
The \texttt{assert(count == 1)} in \texttt{get} fires. Fix: \texttt{while}.
\end{Sol}

\begin{Qn}{Bad schedule}{Single CV with \texttt{while}: everyone asleep}
Now with \texttt{while}, but still one CV. Show how \(T_{c1}, T_{c2}, T_p\)
can all end up asleep. What is the fix?
\end{Qn}
\begin{Sol}
\(T_{c1}\) and \(T_{c2}\) both find count = 0 and sleep. \(T_p\) puts
(count = 1), signals: wakes \(T_{c1}\); loops, finds the buffer full,
sleeps. \(T_{c1}\) re-checks (count = 1), gets (count = 0), and
\textbf{signals the one CV}: this may wake \(T_{c2}\) instead of \(T_p\).
\(T_{c1}\) loops, finds 0, sleeps. \(T_{c2}\) re-checks, finds 0, sleeps.
Now all three are asleep and nobody will signal.

Fix: \textbf{two CVs}. Producers wait on \texttt{empty} and signal
\texttt{fill}; consumers wait on \texttt{fill} and signal \texttt{empty}.
A consumer can then only wake a producer and vice versa. (Broadcast would
also avoid the hang, but two CVs is the clean fix.)
\end{Sol}

\begin{Qn}{Trace}{Multi-slot buffer}
\texttt{MAX = 3}, \texttt{put}/\texttt{get} with \texttt{fill},
\texttt{use}, \texttt{count} as in lecture, initially 0. Operations in
order: put(10), put(11), get, put(12), put(13), put(14), get, get. Track
\texttt{buffer}, \texttt{fill}, \texttt{use}, \texttt{count} and say what
each get returns and where the producer must wait. Assume a producer that
is woken runs (and completes its put) before the next get.
\end{Qn}
\begin{Sol}
\begin{center}
\small
\begin{tabular}{@{}llccccl@{}}
  \toprule
  \textbf{Op} & \textbf{buffer[0..2]} & \textbf{fill} & \textbf{use} & \textbf{count} & \textbf{Returns} & \textbf{Note} \\
  \midrule
  put(10) & 10 -- -- & 1 & 0 & 1 & & \\
  put(11) & 10 11 -- & 2 & 0 & 2 & & \\
  get     & 10 11 -- & 2 & 1 & 1 & 10 & \\
  put(12) & 10 11 12 & 0 & 1 & 2 & & \texttt{fill} wraps \\
  put(13) & 13 11 12 & 1 & 1 & 3 & & full \\
  put(14) & & & & 3 & & \texttt{count == MAX}: producer \textbf{waits} on \texttt{empty} \\
  get     & 13 11 12 & 1 & 2 & 2 & 11 & signals \texttt{empty}; producer wakes, re-checks \\
  (put(14)) & 13 14 12 & 2 & 2 & 3 & & now completes \\
  get     & 13 14 12 & 2 & 0 & 2 & 12 & \texttt{use} wraps \\
  \bottomrule
\end{tabular}
\end{center}
Items come out in FIFO order (10, 11, 12). With one producer and one
consumer, extra slots mean fewer context switches; with several of each,
they can also work concurrently.
\end{Sol}

\begin{Qn}{Concept}{Covering conditions}
In the allocator (\texttt{allocate} waits \texttt{while (bytesLeft <
size)}; \texttt{free} adds bytes and signals), 0 bytes are free.
\(T_a\) calls \texttt{allocate(100)}, \(T_b\) calls \texttt{allocate(10)},
both sleep. \(T_c\) calls \texttt{free(50)}. (a) What can go wrong with
\texttt{signal}? (b) What does \texttt{broadcast} do here, step by step?
(c) What is the cost, and when is needing broadcast a sign of a bug?
\end{Qn}
\begin{Sol}
(a) \texttt{signal} wakes \emph{one} waiter, possibly \(T_a\): it re-checks
\(50 < 100\) and sleeps again, while \(T_b\), who could proceed, stays
asleep. Progress stalls although memory is available.

(b) Both wake. Whichever runs first re-checks: \(T_b\) finds
\(50 \ge 10\), allocates (bytesLeft = 40); \(T_a\) finds \(40 < 100\) (or
\(50<100\)) and sleeps again.

(c) Cost: threads that cannot proceed are woken, run, and go back to sleep
(wasted context switches). If a program only works with broadcast where a
signal should suffice (e.g.\ single-CV producer/consumer), there is
probably a bug (use separate CVs). Broadcast is right when the signaller
cannot know which waiter can proceed, as here.
\end{Sol}

\begin{Qn}{Design}{Strict alternation with a CV}
Two threads must print \texttt{ping} and \texttt{pong} alternately,
starting with \texttt{ping}, 100 times each. Write the code using one mutex,
one CV and a shared variable.
\end{Qn}
\begin{Sol}
\begin{lstlisting}[numbers=none]
int turn = 0;            // 0: ping's turn, 1: pong's turn
mutex_t m; cond_t c;

void *ping(void *arg) {               void *pong(void *arg) {
  for (int i = 0; i < 100; i++) {       for (int i = 0; i < 100; i++) {
    Mutex_lock(&m);                       Mutex_lock(&m);
    while (turn != 0)                     while (turn != 1)
      Cond_wait(&c, &m);                    Cond_wait(&c, &m);
    printf("ping\n");                     printf("pong\n");
    turn = 1;                             turn = 0;
    Cond_signal(&c);                      Cond_signal(&c);
    Mutex_unlock(&m);                     Mutex_unlock(&m);
  }                                     }
}                                     }
\end{lstlisting}
All three rules of thumb: state (\texttt{turn}) kept alongside the CV,
wait/signal with the lock held, \texttt{while} re-checks. One CV is enough
here because there is only ever one waiter, and it is always the thread
whose turn it now is.
\end{Sol}

% ===========================================================================
\section{Lecture 19: Semaphores}
\renewcommand{\lecno}{19}
% ===========================================================================

\begin{Qn}{Concept}{CVs vs.\ semaphores}
(a) What state does each have? (b) Why does join with a semaphore not need
a \texttt{done} variable or a mutex? (c) What does \texttt{sem\_init(\&s,
0, 1)} mean, argument by argument? (d) If a semaphore's value is \(-3\),
what does that tell you (with the lecture's definition)?
\end{Qn}
\begin{Sol}
(a) A CV has no state except its waiting queue; the programmer tracks the
condition. A semaphore has an integer value that determines whether
\texttt{sem\_wait} blocks; the program cannot read it directly.
(b) The value \emph{is} the state: a \texttt{sem\_post} before anybody waits
is remembered (value 1), so the later \texttt{sem\_wait} returns at once;
each operation is atomic, so there is no check-then-sleep gap.
(c) Semaphore \texttt{s}; 0 = shared between threads of the same process;
initial value 1.
(d) Three threads are waiting on it (negative value = number of waiters).
\end{Sol}

\begin{Qn}{Trace}{Semaphore values}
\texttt{s} is initialized to 1. In order: T1 \texttt{sem\_wait}, T2
\texttt{sem\_wait}, T3 \texttt{sem\_wait}, T1 \texttt{sem\_post}, T2
\texttt{sem\_post}, T3 \texttt{sem\_post}. Give the value and each
thread's state after every operation. What does \texttt{s} implement?
\end{Qn}
\begin{Sol}
\begin{center}
\small
\begin{tabular}{@{}lcl@{}}
  \toprule
  \textbf{Operation} & \textbf{value} & \textbf{Effect} \\
  \midrule
  T1 wait & 0  & T1 continues (enters) \\
  T2 wait & \(-1\) & T2 sleeps \\
  T3 wait & \(-2\) & T3 sleeps (2 waiters) \\
  T1 post & \(-1\) & wakes one waiter, say T2 \\
  T2 post & 0  & wakes T3 \\
  T3 post & 1  & nobody waiting \\
  \bottomrule
\end{tabular}
\end{center}
A \textbf{lock} (binary semaphore, init 1): at most one thread between its
wait and post.
\end{Sol}

\begin{Qn}{Fill}{Initial values}
Give the initial value and explain: (a) lock built from a semaphore;
(b) join (parent waits for child); (c) at most 4 threads may use a pool of
4 database connections at once; (d) \texttt{emptyBuffer} and
\texttt{fullBuffer} for an \(N\)-slot buffer. (e) What goes wrong if the
lock semaphore is initialized to 0? to 2? If the join semaphore is
initialized to 1?
\end{Qn}
\begin{Sol}
(a) 1. (b) 0. (c) 4 (counts available resources). (d)
\texttt{emptyBuffer} = \(N\) (the producer can run \(N\) times first),
\texttt{fullBuffer} = 0 (the consumer must wait).

(e) Lock at 0: the very first \texttt{acquire} blocks forever (deadlock).
Lock at 2: two threads can be in the critical section together (no mutual
exclusion). Join at 1: the parent's \texttt{sem\_wait} returns immediately,
even if the child has not finished.
\end{Sol}

\begin{Qn}{Trace}{Join with a semaphore, both orders}
Join uses one semaphore initialized to 0 (parent \texttt{sem\_wait}, child
\texttt{sem\_post}). Trace the value and states when (a) the parent waits
first, (b) the child finishes first. Why did the second order break join
attempt 1 with CVs but not here?
\end{Qn}
\begin{Sol}
(a) Parent wait: \(0 \to -1\), parent sleeps. Child post: \(-1 \to 0\),
wakes parent. Parent's \texttt{sem\_wait} returns (value 0).

(b) Child post: \(0 \to 1\), nobody to wake. Later parent wait:
\(1 \to 0\), not negative, returns immediately.

With CVs, a signal with no waiter is lost, so the parent slept forever. A
semaphore \textbf{remembers} the post in its value.
\end{Sol}

\keepspace{24\baselineskip}
\begin{Qn}{Bad schedule}{Semaphore built from a lock and a CV}
\begin{lstlisting}[numbers=none]
void sem_wait(sem_t *s) {             void sem_post(sem_t *s) {
  lock_acquire(&s->lock);               lock_acquire(&s->lock);
  if (s->value <= 0)          // !!     s->value++;
     cond_wait(&s->cond);               cond_signal(&s->cond);
  s->value--;                           lock_release(&s->lock);
  lock_release(&s->lock);             }
}
\end{lstlisting}
With \texttt{if} instead of \texttt{while}, show how the value can go
negative and two threads can pass a semaphore initialized to 1 used as a
lock.
\end{Qn}
\begin{Sol}
Value 1. T1 \texttt{sem\_wait}: value 0, enters. T2 \texttt{sem\_wait}:
value 0, waits. T1 \texttt{sem\_post}: value 1, signals T2 (T2 is now
ready, not yet running). T3 \texttt{sem\_wait} gets the lock first: value
1 \(> 0\), decrements to 0, enters. T2 finally re-acquires the lock and,
with \texttt{if}, does not re-check: decrements to \(-1\) and enters.
\textbf{T2 and T3 are both inside}. With \texttt{while}, T2 would see 0
and wait again (rule of thumb 3, Mesa semantics).
\end{Sol}

\begin{Qn}{Design}{Ordering threads}
Threads A, B, C each print their letter once. Using semaphores only (give
initial values), force the output order \texttt{C A B}. Then: two threads
must both reach a point before either continues (a rendezvous). Write it,
and explain what happens if each thread waits before it posts.
\end{Qn}
\begin{Sol}
Order C, A, B: \texttt{sA = 0, sB = 0}.
\begin{lstlisting}[numbers=none]
C: print("C"); sem_post(&sA);
A: sem_wait(&sA); print("A"); sem_post(&sB);
B: sem_wait(&sB); print("B");
\end{lstlisting}
Rendezvous: \texttt{aArrived = 0, bArrived = 0}.
\begin{lstlisting}[numbers=none]
A: ...; sem_post(&aArrived); sem_wait(&bArrived); ...
B: ...; sem_post(&bArrived); sem_wait(&aArrived); ...
\end{lstlisting}
Each announces its arrival and then waits for the other; whichever comes
second finds the other's post already remembered. If both \emph{wait first}
(\texttt{sem\_wait} before \texttt{sem\_post}), each waits for a post the
other can only make after its own wait: \textbf{deadlock}.
\end{Sol}

\begin{Qn}{Bad schedule}{Mutex placement in producer/consumer}
Multiple producers/consumers, \(N\) slots, semaphores \texttt{emptyBuffer}
\(= N\), \texttt{fullBuffer} \(= 0\), \texttt{mutex} \(= 1\).
(a) Option \#1 takes \texttt{mutex} \emph{before} \texttt{sem\_wait} on
\texttt{emptyBuffer}/\texttt{fullBuffer}. Give a deadlocking schedule with
semaphore values. (b) Why does option \#2 (mutex inside) not deadlock?
(c) Why is option \#3 (mutex only around \texttt{findempty}/\texttt{findfull})
better? (d) In option \#2, does the order of the two final
\texttt{sem\_post}s matter?
\end{Qn}
\begin{Sol}
(a) Consumer first: \texttt{wait(mutex)} \(1 \to 0\), then
\texttt{wait(fullBuffer)} \(0 \to -1\): sleeps \emph{holding mutex}.
Producer: \texttt{wait(emptyBuffer)} is not reached because it first does
\texttt{wait(mutex)} \(0 \to -1\): sleeps. Nobody can post: deadlock.

(b) A thread takes \texttt{mutex} only after its resource wait succeeded,
so it never sleeps while holding \texttt{mutex}; the critical section
always finishes and releases it.

(c) In \#2 the whole \texttt{Fill}/\texttt{Use} is inside \texttt{mutex}:
only one thread at a time works on \emph{any} slot. In \#3 only finding
(and claiming) a slot is serialized; filling/using different slots proceeds
concurrently.

(d) No: posts never block, so either order releases both semaphores.
(Only the order of the \emph{waits} matters.)
\end{Sol}

\begin{Qn}{Concept}{Private or shared indices}
Why does the circular-buffer code (\texttt{i = (i+1)\%N} in the producer,
\texttt{j} in the consumer) work with one producer and one consumer but fail
with several producers? Consider \texttt{i} private and \texttt{i} shared.
\end{Qn}
\begin{Sol}
One producer and one consumer: only the producer touches \texttt{i} and only
the consumer touches \texttt{j}; the semaphores guarantee they never work
on the same slot, so no lock is needed.

Several producers, \texttt{i} \textbf{private}: each starts at 0 and they
fill the \emph{same} slots, overwriting each other. \texttt{i}
\textbf{shared}: two producers can read the same \texttt{i} before either
increments it (a race, like the shared counter). Either way each producer
must claim a unique slot under mutual exclusion.
\end{Sol}

\keepspace{22\baselineskip}
\begin{Qn}{Trace}{Reader/writer lock}
\texttt{lock} = 1, \texttt{writelock} = 1, \texttt{readers} = 0.
Events: R1 acquire-read, R2 acquire-read, W acquire-write, R3
acquire-read, R1 release, R2 release, R3 release. Give \texttt{readers},
\texttt{writelock} and who is inside/waiting after each event. What
fairness problem does this show?
\end{Qn}
\begin{Sol}
\begin{center}
\small
\begin{tabular}{@{}lccl@{}}
  \toprule
  \textbf{Event} & \textbf{readers} & \textbf{writelock} & \textbf{State} \\
  \midrule
  R1 acquire & 1 & 0 & first reader takes writelock; R1 reading \\
  R2 acquire & 2 & 0 & R1, R2 reading \\
  W acquire  & 2 & \(-1\) & W waits on writelock \\
  R3 acquire & 3 & \(-1\) & R3 \textbf{enters ahead of W} \\
  R1 release & 2 & \(-1\) & \\
  R2 release & 1 & \(-1\) & \\
  R3 release & 0 & 0 & last reader posts writelock: \textbf{W enters} \\
  \bottomrule
\end{tabular}
\end{center}
(\texttt{lock} goes 1 \(\to\) 0 \(\to\) 1 around each reader's update of
\texttt{readers}.) As long as new readers keep arriving before the count
reaches 0, the writer \textbf{starves}. Fixing it requires stopping new
readers once a writer is waiting (the open question on the slide).
\end{Sol}

\begin{Qn}{Concept}{Dining philosophers}
(a) Show the deadlock in the basic solution
(\texttt{wait(left); wait(right)}). (b) Why does making philosopher 4
take right first remove it? (c) Lefty/righty: if every philosopher is a
righty, is deadlock possible? What if there is at least one of each?
(d) What is the maximum number of philosophers that can eat at once with 5
forks?
\end{Qn}
\begin{Sol}
(a) Each \(P_i\) takes \(f_i\) (its left) before anyone takes a right
fork. Then \(P_i\) waits for \(f_{i+1}\), held by \(P_{i+1}\), \ldots,
\(P_4\) waits for \(f_0\), held by \(P_0\): a \textbf{cycle}; everyone
waits forever.

(b) \(P_4\) asks for \(f_0\) first. If \(P_0\) holds \(f_0\), \(P_4\) waits
\emph{without} holding \(f_4\), so \(P_3\) can take \(f_4\), eat and release;
the chain unwinds. The circular wait is impossible.

(c) All righties: same as all lefties by symmetry, so deadlock is
possible. At least one lefty and one righty: the acquisition order is no
longer uniform around the circle, so the cycle cannot form (no deadlock;
starvation freedom is a separate question).

(d) \textbf{2}: each eater uses 2 of the 5 forks, and
\(\lfloor 5/2 \rfloor = 2\) (e.g.\ \(P_0\) and \(P_2\)).
\end{Sol}

\begin{Qn}{Design}{A one-doctor clinic (past-exam style)}
One doctor, unlimited waiting room, one shared buffer. A patient calls
\texttt{consult()} (writes symptoms to the buffer) and later
\texttt{readTreatment()}; the doctor loops calling \texttt{treat()} (reads
symptoms, writes treatment). Using only semaphores: patients must not
overwrite each other, the doctor must only read new symptoms, and a patient
must only read a finished treatment.
\end{Qn}
\begin{Sol}
\texttt{office = 1} (one patient uses the buffer), \texttt{symptoms = 0}
(new request ready), \texttt{treated = 0} (treatment ready).
\begin{lstlisting}[numbers=none]
Patient:                          Doctor:
  sem_wait(&office);                while (1) {
  consult();                          sem_wait(&symptoms);
  sem_post(&symptoms);                treat();
  sem_wait(&treated);                 sem_post(&treated);
  readTreatment();                  }
  sem_post(&office);
\end{lstlisting}
\texttt{office} works as a lock held for the whole request--response
exchange; the two 0-initialized semaphores are join-style ordering
(``symptoms written before treat'', ``treatment written before read''). The
waiting room is the queue of \texttt{office}.
\end{Sol}

\begin{Qn}{Concept}{Equivalence}
``Semaphores are equally powerful to locks + CVs.'' (a) What has to be
shown? (b) Which direction did the lecture call really hard? (c) Write
\texttt{sem\_wait}/\texttt{sem\_post} using a lock and a CV.
\end{Qn}
\begin{Sol}
(a) Each can be built from the other: locks from semaphores, CVs from
semaphores, and semaphores from locks + CVs (convenience is irrelevant).
(b) Building \textbf{CVs from semaphores} (possible, but hard to get right).
(c)
\begin{lstlisting}[numbers=none]
void sem_wait(sem_t *s) {               void sem_post(sem_t *s) {
  lock_acquire(&s->lock);                 lock_acquire(&s->lock);
  while (s->value <= 0)                   s->value++;
    cond_wait(&s->cond, &s->lock);        cond_signal(&s->cond);
  s->value--;                             lock_release(&s->lock);
  lock_release(&s->lock);               }
}
\end{lstlisting}
(Here the value never goes below 0, so it does not count waiters.)
\end{Sol}

% ===========================================================================
\section{Mixed / exam-style}
\renewcommand{\lecno}{M}
% ===========================================================================

\keepspace{17\baselineskip}
\begin{Qn}{T/F}{Across lectures}
\begin{enumerate}[(a)]
  \item A thread blocked on a semaphore and a process blocked on a page
        fault are both removed from the ready queue.
  \item A thread may go to sleep (e.g.\ wait for I/O) while holding a lock
        without violating mutual exclusion.
  \item If \(T_2\) writes a variable without taking the mutex that \(T_1\)
        uses to protect it, the hardware generates a trap.
  \item A context switch can happen only after a timer interrupt.
  \item Spinlocks are a good choice on a uniprocessor.
\end{enumerate}
\end{Qn}
\begin{Sol}
(a) \textbf{True}: both wait for an event (post / disk I/O), so the
scheduler runs other ready work.
(b) \textbf{True}: other threads still cannot enter (it may be slow or
deadlock-prone, but exclusion holds).
(c) \textbf{False}: a mutex is a software convention; memory accesses carry
no lock information. It is a data race, not a trap.
(d) \textbf{False}: also when a thread blocks (I/O, page fault,
\texttt{sem\_wait}, \texttt{cond\_wait}, \texttt{park}) or yields.
(e) \textbf{False}: the holder cannot run while the waiter spins; waste of a
whole slice per waiter.
\end{Sol}

\begin{Qn}{Design}{Bounded buffer from scratch}
Write a correct bounded buffer of size \(N\) for multiple producers and
consumers (a) with locks + CVs and (b) with semaphores. For each, point to
the line that would break if changed (\texttt{while} \(\to\) \texttt{if};
mutex order).
\end{Qn}
\begin{Sol}
(a) Mutex \texttt{m}, CVs \texttt{empty}, \texttt{fill}, \texttt{count}:
producer \texttt{lock; while (count == N) wait(empty); put; signal(fill);
unlock}; consumer \texttt{lock; while (count == 0) wait(fill); get;
signal(empty); unlock}. \texttt{while} \(\to\) \texttt{if}: a woken consumer
may find the item taken by another consumer (bad data).

(b) \texttt{emptyBuffer = N, fullBuffer = 0, mutex = 1}: producer
\texttt{wait(emptyBuffer); wait(mutex); myi = findempty; post(mutex);
Fill; post(fullBuffer)}; consumer symmetric. Taking \texttt{mutex} before
\texttt{emptyBuffer}/\texttt{fullBuffer}: a thread can sleep holding
\texttt{mutex} $\Rightarrow$ deadlock.
\end{Sol}

\end{document}
